\begin{figure}[ht]
\centering
\begin{tikzpicture}[x=1cm, y=1cm, font=\small]

% Four panels showing the distribution of sample means tightening with n
\foreach \n/\xshift/\sd in {1/0/2.4, 5/4.5/1.07, 10/9/0.76, 100/13.5/0.24} {
  \begin{scope}[shift={(\xshift,0)}]
    \node[anchor=south] at (1.5, 3.6) {$n = \n$};
    \draw[->, thick] (-0.2, 0) -- (3.2, 0);
    \draw[->, thick] (0, -0.1) -- (0, 3.4);
    \draw[thick, engblue, dashed] (1.5, 0) -- (1.5, 3.0);
    % Gaussian-ish bell centered at 1.5
    \draw[thick, engred] plot[domain=0.05:2.95, samples=40]
      (\x, {3.0*exp(-((\x-1.5)*(\x-1.5))/(2*\sd*\sd*0.18))});
    \node[anchor=north, font=\scriptsize] at (1.5, -0.2) {sample mean};
  \end{scope}
}

\node[anchor=north, font=\small] at (8, -1.1) {
  As $n$ grows, the distribution of the sample mean tightens as
  $\sigma/\sqrt{n}$ and becomes Gaussian (CLT).
};

\end{tikzpicture}
\caption[Central limit theorem convergence with sample size]{The
distribution of the sample mean tightens as $1/\sqrt{n}$ and tends
to a Gaussian shape regardless of the parent distribution (here
shown as the bell curve at $n = 1$ already approximately Gaussian
for illustration; for a uniform parent the $n = 1$ panel would be
flat-topped). The reduction in spread is the random-error component
of measurement uncertainty.}
\label{fig:vol01:ch04:clt-convergence}
\end{figure}
